\uuid{CXUp}
\chapitre{Probabilité continue}
\niveau{L2}
\module{Probabilité}
\sousChapitre{Loi normale}
\titre{Calculs avec une loi normale}
\theme{loi normale}

\auteur{}
\datecreate{2022-08-25}
\organisation{AMSCC}
\difficulte{}
\contenu{
\texte{$\mathcal N(\mu,\sigma)$ désigne la loi normale de moyenne $\mu$ et d’écart-type $\sigma$.}
\texte{On suppose que la durée annuelle de travail d'une personne non salariée suit une loi normale d'espérance 2600 heures et d'écart-type 200 heures.}
\question{Quelle est la probabilité de travailler moins de 3000 heures ? Entre 2500 et 2700 heures ? Utiliser la table puis vérifier avec un tableur.}
\indication{Réduire la variable : $Z=(X-2600)/200$.}
\reponse{Avec $X\sim\mathcal N(2600,200)$, $Z\sim\mathcal N(0,1)$.
\[\mathbb P(X<3000)=\Phi(2)\simeq0{,}97725,\qquad
\mathbb P(2500\leq X\leq2700)=\Phi(0{,}5)-\Phi(-0{,}5)\simeq0{,}38292.\]
Avec Excel en français : \texttt{=LOI.NORMALE.N(3000;2600;200;VRAI)} ; pour le second résultat, soustraire les valeurs cumulées en 2700 et 2500.}
}
